% Standalone research note for the hybrid-local-solver project.

\documentclass[11pt]{article}

% Common class-neutral preamble for standalone notes under manuscript/notes/.
% Note entry points all live exactly two directories below manuscript/.

\usepackage[margin=1in]{geometry}
\usepackage{amsmath,amssymb,amsthm,mathtools,bm}
\usepackage{algorithm}
\usepackage{algorithmic}
\usepackage{booktabs,tabularx,multirow,array,longtable}
\usepackage{enumitem}
\usepackage{microtype}
\usepackage[numbers,sort&compress]{natbib}
\usepackage{xcolor}
\usepackage[colorlinks=true,citecolor=blue,linkcolor=black,urlcolor=blue]{hyperref}
\usepackage[capitalize,nameinlink,noabbrev]{cleveref}

\newtheorem{theorem}{Theorem}[section]
\newtheorem{lemma}[theorem]{Lemma}
\newtheorem{proposition}[theorem]{Proposition}
\newtheorem{corollary}[theorem]{Corollary}
\newtheorem{conjecture}[theorem]{Conjecture}
\theoremstyle{definition}
\newtheorem{definition}[theorem]{Definition}
\newtheorem{assumption}[theorem]{Assumption}
\newtheorem{problem}[theorem]{Problem}
\newtheorem{observation}[theorem]{Observation}
\theoremstyle{remark}
\newtheorem{remark}[theorem]{Remark}
\theoremstyle{plain}

% Canonical mathematical notation for the active manuscript.
%
% This file preserves the recurring notation style used in the author's
% NeurIPS 2024 and 2025 papers while excluding unrelated tensor and
% machine-learning boilerplate from those archived command collections.
%
% Prefix convention:
%   \v...  bold vectors
%   \m...  bold matrices
%   \g...  calligraphic sets, graphs, and families
%   \s...  blackboard-bold spaces and number systems

\newcommand{\method}{\textsc{HybridLocalSolver}}

% Font and scalar shorthands retained from the archived manuscripts.
\newcommand{\mc}{\mathcal}
\newcommand{\R}{\mathbb{R}}
\newcommand{\E}{\mathbb{E}}
\newcommand{\G}{\mathbb{G}}
\newcommand{\N}{\mathcal{N}}
% Accuracy namespaces are semantic: do not add a bare \eps alias.
% Degree-normalized residual tolerance of classical APPR is kept distinct from
% the objective-gap and proximal fixed-point tolerances of the RPPR sections.
\newcommand{\epsappr}{\varepsilon_{\mathrm{appr}}}
\newcommand{\epsobj}{\varepsilon_{\mathrm{obj}}}
\newcommand{\epspg}{\varepsilon_{\mathrm{pg}}}
\newcommand{\epsppr}{\varepsilon_{\mathrm{ppr}}}
\newcommand{\epsin}{\varepsilon_{\mathrm{in}}}
\newcommand{\epskkt}{\varepsilon_{\mathrm{kkt}}}
\newcommand{\epsburn}{\varepsilon_{\mathrm{burn}}}
\newcommand{\epssol}{\varepsilon_{\mathrm{sol}}}
\newcommand{\epspprinst}[1]{\varepsilon_{\mathrm{ppr},#1}}
\newcommand{\trans}{\mathsf{T}}
\newcommand{\grad}{\nabla}

% Delimiters and generic helpers.
% Norms and inner products use fixed-size delimiters by default.
% Use an optional size (e.g. \norm[\big]{...}) or * for automatic sizing.
\DeclarePairedDelimiter{\norm}{\lVert}{\rVert}
\newcommand{\abs}[1]{\left\lvert #1 \right\rvert}
\DeclarePairedDelimiterX{\inner}[2]{\langle}{\rangle}{#1, #2}
\newcommand{\ip}[2]{\inner{#1}{#2}}
\newcommand{\card}[1]{\left\lvert #1 \right\rvert}
\newcommand{\set}[1]{\left\{#1\right\}}
\newcommand{\ceil}[1]{\left\lceil #1 \right\rceil}
\newcommand{\floor}[1]{\left\lfloor #1 \right\rfloor}
\newcommand{\comp}[1]{\overline{#1}}
\newcommand{\conj}[1]{\overline{#1}}
\newcommand{\mean}[1]{\overline{#1}}
\newcommand{\bigO}{\mathcal{O}}
\newcommand{\softO}{\widetilde{\mathcal{O}}}
\newcommand{\softOmega}{\widetilde{\Omega}}
\newcommand{\pos}[1]{\left[#1\right]_{+}}

% Bold lowercase vectors.
\newcommand{\va}{\bm{a}}
\newcommand{\vb}{\bm{b}}
\newcommand{\vc}{\bm{c}}
\newcommand{\vd}{\bm{d}}
\newcommand{\ve}{\bm{e}}
\newcommand{\vf}{\bm{f}}
\newcommand{\vg}{\bm{g}}
\newcommand{\vh}{\bm{h}}
\newcommand{\vi}{\bm{i}}
\newcommand{\vj}{\bm{j}}
\newcommand{\vk}{\bm{k}}
\newcommand{\vl}{\bm{l}}
\newcommand{\vm}{\bm{m}}
\newcommand{\vn}{\bm{n}}
\newcommand{\vo}{\bm{o}}
\newcommand{\vp}{\bm{p}}
\newcommand{\vq}{\bm{q}}
\newcommand{\vr}{\bm{r}}
\newcommand{\vs}{\bm{s}}
\newcommand{\vt}{\bm{t}}
\newcommand{\vu}{\bm{u}}
\newcommand{\vv}{\bm{v}}
\newcommand{\vw}{\bm{w}}
\newcommand{\vx}{\bm{x}}
\newcommand{\vy}{\bm{y}}
\newcommand{\vz}{\bm{z}}

% Bold Greek vectors and special vectors.
\newcommand{\valpha}{\bm{\alpha}}
\newcommand{\vbeta}{\bm{\beta}}
\newcommand{\vdelta}{\bm{\delta}}
\newcommand{\vepsilon}{\bm{\epsilon}}
\newcommand{\veta}{\bm{\eta}}
\newcommand{\vgamma}{\bm{\gamma}}
\newcommand{\vlambda}{\bm{\lambda}}
\newcommand{\vmu}{\bm{\mu}}
\newcommand{\vomega}{\bm{\omega}}
\newcommand{\vphi}{\bm{\phi}}
\newcommand{\vpi}{\bm{\pi}}
\newcommand{\vpsi}{\bm{\psi}}
\newcommand{\vrho}{\bm{\rho}}
\newcommand{\vsigma}{\bm{\sigma}}
\newcommand{\vtheta}{\bm{\theta}}
\newcommand{\vxi}{\bm{\xi}}
\newcommand{\vell}{\bm{\ell}}
\newcommand{\vDelta}{\bm{\Delta}}
\newcommand{\vGamma}{\bm{\Gamma}}
\newcommand{\vLambda}{\bm{\Lambda}}
\newcommand{\vPhi}{\bm{\Phi}}
\newcommand{\vSigma}{\bm{\Sigma}}
\newcommand{\vzero}{\bm{0}}
\newcommand{\vone}{\bm{1}}
\newcommand{\one}{\mathbf{1}}
\newcommand{\zero}{\vzero}
\newcommand{\eunit}[1]{\bm{e}_{#1}}
\newcommand{\vDeltax}{\bm{\Delta}\bm{x}}

% Bold uppercase matrices.
\newcommand{\mA}{\bm{A}}
\newcommand{\mB}{\bm{B}}
\newcommand{\mC}{\bm{C}}
\newcommand{\mD}{\bm{D}}
\newcommand{\mE}{\bm{E}}
\newcommand{\mF}{\bm{F}}
\newcommand{\mG}{\bm{G}}
\newcommand{\mH}{\bm{H}}
\newcommand{\mI}{\bm{I}}
\newcommand{\mJ}{\bm{J}}
\newcommand{\mK}{\bm{K}}
\newcommand{\mL}{\bm{L}}
\newcommand{\mM}{\bm{M}}
\newcommand{\mN}{\bm{N}}
\newcommand{\mO}{\bm{O}}
\newcommand{\mP}{\bm{P}}
\newcommand{\mQ}{\bm{Q}}
\newcommand{\mR}{\bm{R}}
\newcommand{\mS}{\bm{S}}
\newcommand{\mT}{\bm{T}}
\newcommand{\mU}{\bm{U}}
\newcommand{\mV}{\bm{V}}
\newcommand{\mW}{\bm{W}}
\newcommand{\mX}{\bm{X}}
\newcommand{\mY}{\bm{Y}}
\newcommand{\mZ}{\bm{Z}}

% Bold Greek matrices retained from the graph papers.
\newcommand{\mBeta}{\bm{\beta}}
\newcommand{\mDelta}{\bm{\Delta}}
\newcommand{\mGamma}{\bm{\Gamma}}
\newcommand{\mLambda}{\bm{\Lambda}}
\newcommand{\mOmega}{\bm{\Omega}}
\newcommand{\mPhi}{\bm{\Phi}}
\newcommand{\mPi}{\bm{\Pi}}
\newcommand{\mSigma}{\bm{\Sigma}}
\newcommand{\mTheta}{\bm{\Theta}}

% Calligraphic symbols.
\newcommand{\gA}{\mathcal{A}}
\newcommand{\gB}{\mathcal{B}}
\newcommand{\gC}{\mathcal{C}}
\newcommand{\gD}{\mathcal{D}}
\newcommand{\gE}{\mathcal{E}}
\newcommand{\gF}{\mathcal{F}}
\newcommand{\gG}{\mathcal{G}}
\newcommand{\gH}{\mathcal{H}}
\newcommand{\gI}{\mathcal{I}}
\newcommand{\gJ}{\mathcal{J}}
\newcommand{\gK}{\mathcal{K}}
\newcommand{\gL}{\mathcal{L}}
\newcommand{\gM}{\mathcal{M}}
\newcommand{\gN}{\mathcal{N}}
\newcommand{\gO}{\mathcal{O}}
\newcommand{\gP}{\mathcal{P}}
\newcommand{\gQ}{\mathcal{Q}}
\newcommand{\gR}{\mathcal{R}}
\newcommand{\gS}{\mathcal{S}}
\newcommand{\gT}{\mathcal{T}}
\newcommand{\gU}{\mathcal{U}}
\newcommand{\gV}{\mathcal{V}}
\newcommand{\gW}{\mathcal{W}}
\newcommand{\gX}{\mathcal{X}}
\newcommand{\gY}{\mathcal{Y}}
\newcommand{\gZ}{\mathcal{Z}}

% Blackboard-bold symbols.
\newcommand{\sA}{\mathbb{A}}
\newcommand{\sB}{\mathbb{B}}
\newcommand{\sC}{\mathbb{C}}
\newcommand{\sD}{\mathbb{D}}
\newcommand{\sE}{\mathbb{E}}
\newcommand{\sF}{\mathbb{F}}
\newcommand{\sG}{\mathbb{G}}
\newcommand{\sH}{\mathbb{H}}
\newcommand{\sI}{\mathbb{I}}
\newcommand{\sJ}{\mathbb{J}}
\newcommand{\sK}{\mathbb{K}}
\newcommand{\sL}{\mathbb{L}}
\newcommand{\sM}{\mathbb{M}}
\newcommand{\sN}{\mathbb{N}}
\newcommand{\sO}{\mathbb{O}}
\newcommand{\sP}{\mathbb{P}}
\newcommand{\sQ}{\mathbb{Q}}
\newcommand{\sR}{\mathbb{R}}
\newcommand{\sS}{\mathbb{S}}
\newcommand{\sT}{\mathbb{T}}
\newcommand{\sU}{\mathbb{U}}
\newcommand{\sV}{\mathbb{V}}
\newcommand{\sW}{\mathbb{W}}
\newcommand{\sX}{\mathbb{X}}
\newcommand{\sY}{\mathbb{Y}}
\newcommand{\sZ}{\mathbb{Z}}

% Frequently used operators.
\DeclareMathOperator{\Cov}{Cov}
\DeclareMathOperator{\Var}{Var}
\DeclareMathOperator{\diag}{diag}
\DeclareMathOperator{\nei}{Nei}
\DeclareMathOperator{\nnz}{nnz}
\DeclareMathOperator{\pr}{pr}
\DeclareMathOperator{\prox}{prox}
\DeclareMathOperator{\push}{push}
\DeclareMathOperator{\res}{res}
\DeclareMathOperator{\rel}{Rel}
\DeclareMathOperator{\relu}{ReLU}
\DeclareMathOperator{\sign}{sign}
\DeclareMathOperator{\supp}{supp}
\DeclareMathOperator{\tr}{tr}
\DeclareMathOperator{\Tr}{Tr}
\DeclareMathOperator{\vol}{vol}
\DeclareMathOperator{\work}{work}
\DeclareMathOperator{\Work}{Work}
\DeclareMathOperator*{\argmax}{arg\,max}
\DeclareMathOperator*{\argmin}{arg\,min}

% Project-wide names and claim-status labels used by standalone research notes.
% Mathematical notation belongs in math_commands.tex.

% Algorithms and solver families.
\newcommand{\AESP}{\textsc{AESP}}
\newcommand{\APPR}{\textsc{APPR}}
\newcommand{\ASPR}{\textsc{ASPR}}
\newcommand{\FISTA}{\textsc{FISTA}}
\newcommand{\ISTA}{\textsc{ISTA}}
\newcommand{\LOCAPPR}{\textsc{LocAPPR}}
\newcommand{\LOCGD}{\textsc{LocGD}}
\newcommand{\LOCSOR}{\textsc{LocSOR}}
\newcommand{\LOCCH}{\textsc{LocCH}}
\newcommand{\LOCHB}{\textsc{LocHB}}
\newcommand{\LocProxGS}{\textsc{LocProxGS}}
\newcommand{\Hybrid}{\textsc{AESP--LocSOR}}

% Claim classifications required by the repository research-note protocol.
\newcommand{\statussource}{\textbf{Source result}}
\newcommand{\statusproved}{\textbf{Proved here}}
\newcommand{\statusconditional}{\textbf{Conditional}}
\newcommand{\statusconjecture}{\textbf{Open / conjectural}}
\newcommand{\statusopen}{\textbf{Open}}
\newcommand{\statusempirical}{\textbf{Empirical}}
\newcommand{\statusobservation}{\textbf{Empirical observation}}
\newcommand{\statusrefuted}{\textbf{Refuted route}}

% Compatibility names used by the older complete-note presentation layer.
\newcommand{\Status}[2]{\textbf{#1:} #2}
\newcommand{\SourceStatus}{\textnormal{\textsc{Source result}}}
\newcommand{\ProvedStatus}{\textnormal{\textsc{Proved here}}}
\newcommand{\ConditionalStatus}{\textnormal{\textsc{Conditional}}}
\newcommand{\ComputedStatus}{\textnormal{\textsc{Computed}}}
\newcommand{\RefutedStatus}{\textnormal{\textsc{Refuted route}}}
\newcommand{\OpenStatus}{\textnormal{\textsc{Open}}}



\title{A Conditional Modal Route to a Terminal Logarithmic Block on Short Paths}
\author{Baojian Zhou\\
\small Working research note for \texttt{hybrid-local-solver}}
\date{August 2026}

\begin{document}
\maketitle

\begin{abstract}
We revisit one quarantined terminal-face candidate for a transported-center
accelerated RPPR execution on the endpoint path $P_m$.  Its parameters are
$q=1/(16m)$, $\alpha=q^2$, and $\rho=\tau=q/5$.  On a projection-inactive
full face, the normalized residual is an exactly solvable damped cosine wave.
We give the literal degree-weighted cosine coefficients, keep the stationary
mode separate from residual range, and prove a phase-uniform
anti-cancellation bound over a band of even modes.

The measured full-face entry is close to the binomial packet
\[
 g_j=-\frac{q^2}{2}(1-q)^m2^{-m}\binom mj.
\]
This packet has about $\sqrt m$ even cosine coefficients of order $q^3$.
We prove that a quantified entry approximation on only
$H_m=\Theta(\sqrt{m/\log m})$ modes, together with uniform projection
inactivity and an unclipped global safe-envelope subtraction, forces at least
$\frac18q^{-1}\log(1/q)$ terminal steps for all sufficiently large $m$.
The two hypotheses are explicit missing lemmas, not conclusions.  A
deterministic floating screen through $m=2048$ supports both, but establishes
no asymptotic result.  Thus this note supplies a rigorous conditional route
and a falsification interface; it does not prove that the logarithm is
necessary, and it gives no lower bound for other recurrences or implicit
response algorithms.
\end{abstract}

% Canonical source-aligned PageRank/RPPR reference formulation.
%
% This fragment is intentionally heading-free at the section level so that the
% active paper and every standalone note can input it. It is a manuscript
% reference convention, not an implementation-wide residual decision.
% Primary source: Fountoulakis and Martinez-Rubio, arXiv:2602.21138v2,
% PDF p. 3, Section 3 and equation (RPPR).

\subsection{Common source-aligned PageRank and RPPR model}
\label{subsec:shared-source-aligned-problem}

All documents in this manuscript workspace use the following reference
language. It follows the plain italic notation of the comparison paper:
\(x,s,A,D,Q,I\) denote column vectors or matrices as determined by their
displayed domains. This is the scoped exception to the project's usual
bold-vector and bold-matrix typography. A note may impose a stronger
assumption after this block, but it must not redefine these objects.

Let \(G=(V,E)\) be a finite, simple, undirected, unweighted graph without
isolated vertices, let \(V=[n]:=\{1,\ldots,n\}\), and let
\(A\in\{0,1\}^{n\times n}\) be its symmetric adjacency matrix. For
\(i\in V\), write
\[
    \mathcal N(i):=\{j\in V:\{i,j\}\in E\},
    \qquad
    d_i:=|\mathcal N(i)|,
    \qquad
    D:=\diag(d_1,\ldots,d_n).
\]
For \(S\subseteq V\), define
\[
    \mathcal N(S):=\bigcup_{i\in S}\mathcal N(i),
    \qquad
    \partial S:=\mathcal N(S)\setminus S,
    \qquad
    \operatorname{Ext}(S):=V\setminus(S\cup\partial S),
\]
\begin{equation}
    \vol(S):=\sum_{i\in S}d_i.
    \label{eq:shared-volume}
\end{equation}
The support of a vector is
\(\supp(x):=\{i\in[n]:x_i\neq0\}\). Matrix inequalities use the Loewner
order, and all unqualified vector inequalities are coordinatewise.

The seed is a column distribution
\begin{equation}
    s\in\R_+^n,
    \qquad
    \inner{\one}{s}=1.
    \label{eq:shared-seed}
\end{equation}
The single-seed specialization is always written \(s=e_v\), where \(v\in V\)
is the seed vertex. Thus \(s\) is never used as a vertex label. For
\(\alpha\in(0,1]\), define
\begin{equation}
\begin{aligned}
    \mathcal L&:=I-D^{-1/2}AD^{-1/2},\\
    Q&:=\alpha I+\frac{1-\alpha}{2}\mathcal L
      =\frac{1+\alpha}{2}I
       -\frac{1-\alpha}{2}D^{-1/2}AD^{-1/2},\\
    b&:=\alpha D^{-1/2}s.
\end{aligned}
    \label{eq:shared-pagerank-matrices}
\end{equation}
Since \(0\preceq\mathcal L\preceq2I\),
\(\alpha I\preceq Q\preceq I\). The shared smooth PageRank quadratic is
\begin{equation}
    f(x):=\frac12\inner{x}{Qx}-\inner{b}{x},
    \qquad
    \nabla f(x)=Qx-b.
    \label{eq:shared-pagerank-objective}
\end{equation}
Its unique unregularized minimizer and associated lazy personalized PageRank
vector are
\begin{equation}
    x^0:=Q^{-1}b,
    \qquad
    \pi:=D^{1/2}x^0.
    \label{eq:shared-ppr-solution}
\end{equation}

For a regularization parameter \(\rho>0\), reserve \(g_\rho\) for the
nonsmooth weighted \(\ell_1\) term and define
\begin{equation}
    g_\rho(x):=\alpha\rho\norm{D^{1/2}x}_1,
    \qquad
    F_\rho(x):=f(x)+g_\rho(x),
    \qquad
    x^\star(\rho):=\argmin_{x\in\R^n}F_\rho(x).
    \label{eq:shared-rppr-objective}
\end{equation}
When \(\rho\) is fixed, \(x^\star\) abbreviates \(x^\star(\rho)\), never the
unregularized point \(x^0\). Write
\[
    S^\star(\rho):=\supp(x^\star(\rho)),
    \qquad
    I^\star(\rho):=[n]\setminus S^\star(\rho).
\]
The source records \(x^\star(\rho)\geq0\),
\(\vol(S^\star(\rho))\leq1/\rho\), and the coordinatewise conditions
\begin{equation}
    \nabla_i f(x^\star(\rho))
    \in
    \begin{cases}
        \{-\alpha\rho\sqrt{d_i}\}, & x_i^\star(\rho)>0,\\
        [-\alpha\rho\sqrt{d_i},0], & x_i^\star(\rho)=0.
    \end{cases}
    \label{eq:shared-rppr-kkt}
\end{equation}

The common computational unit is one repeated adjacency-list scan: scanning
coordinate \(i\) costs \(d_i\), and scanning a set \(S\) costs \(\vol(S)\).
Iteration count and wall-clock time do not replace this degree-weighted work
measure.

Finally, accuracy symbols remain distinct. The APPR threshold is
\(\epsappr\); \(\epsobj\) denotes an objective-gap target;
\(\epspg\) denotes the source experiment's proximal fixed-point tolerance; and
\(\epsppr\) may denote a document-scoped degree-normalized PPR error target.
No equality or conversion among these quantities is assumed. In particular,
this shared model does not resolve the repository-wide residual decision.


\section{Named execution and claim boundary}
\label{sec:terminal-modal-scope}

\paragraph{Algorithmic scope.}
Let $P_m=(v_0,\ldots,v_m)$ be the complete endpoint path with ambient degrees
$d_0=d_m=1$ and $d_j=2$ for $0<j<m$.  Put
\begin{equation}
 q=\frac1{16m},\qquad \alpha=q^2,\qquad \rho=\tau=\frac q5.
 \label{eq:terminal-modal-parameters}
\end{equation}
The named execution starts at zero on $\{v_0\}$, applies one projected
transported-center step per visited prefix, uses the complete all-violations
gate to admit the next singleton, and transports the estimate center by the
exact restricted-optimum displacement.  After admitting $v_m$, it stays on
the full path until the first one-sided certificate and makes one terminal
external return.  These are the exact semantics imported from the sibling
\texttt{volume\_gated\_acceleration} note.

At full-face entry the newly admitted coordinate is still zero-padded.  The
safe-envelope subtraction cannot make it positive, and the same strict
outside-row violation that triggered its admission remains at that row.
Thus entry itself is not an instantaneous terminal certificate; the first
terminal step genuinely occurs after the full face has been entered.

In normalized path coordinates $p=D^{-1/2}x$, put
\begin{equation}
 \widehat Q=D^{-1/2}QD^{1/2}
 =aI-\eta P_{\rm row},
 \quad
 a=\frac{1+q^2}{2},\quad
 \eta=\frac{1-q^2}{2},\quad
 P_{\rm row}=D^{-1}A.
 \label{eq:terminal-modal-operator}
\end{equation}
The normalized affine load has
$\bar c_0=\alpha(1-\rho)$ and
$\bar c_j=-\alpha\rho$ for $j>0$.  Write
\begin{equation}
 r(p):=\widehat Qp-\bar c.
 \label{eq:terminal-modal-residual}
\end{equation}
On a fixed face the step is
\begin{equation}
 y_k=\frac{p_k+qv_k}{1+q},\qquad
 p_{k+1}=\left[y_k-r(y_k)\right]_+,qquad
 v_{k+1}=p_{k+1}+\frac{1-q}{q}(p_{k+1}-p_k).
 \label{eq:terminal-modal-step}
\end{equation}

\subsection*{Claim classification}
The spectral identities, range implication, packet transform,
anti-cancellation lemma, and conditional theorem below are \emph{proved
here}.  The actual entry comparison in
\cref{lem:terminal-modal-packet} and regime preservation in
\cref{lem:terminal-modal-regime} are \emph{open}.  The table in
\cref{sec:terminal-modal-screen} is \emph{measured}.  In particular, no
unconditional terminal logarithmic block or algorithm lower bound is claimed.

\section{Exact full-face modes and the global correction}
\label{sec:terminal-modal-exact}

Let $r_k^x=r(p_k)$ and $r_k^v=r(v_k)$ after entry to the full face, with entry
indexed by $k=0$.  Whenever the projection in
\cref{eq:terminal-modal-step} is inactive, the sibling note proves
\begin{align}
 r_{k+1}^x&=\frac{B}{1+q}(r_k^x+q r_k^v),
 &q r_{k+1}^v&=r_{k+1}^x-(1-q)r_k^x,
 \label{eq:terminal-modal-two-state}\\
 r_{k+1}^x&=B\left(\frac2{1+q}r_k^x
 -\frac{1-q}{1+q}r_{k-1}^x\right),
 &B&=I-\widehat Q.
 \label{eq:terminal-modal-second-order}
\end{align}

For $0\le h\le m$, define
\begin{equation}
 u_h(j)=\cos\frac{h\pi j}{m},\qquad 0\le j\le m.
 \label{eq:terminal-modal-cosines}
\end{equation}
Use the reversible inner product
$\langle z,w\rangle_D=\sum_{j=0}^m d_jz_jw_j$.

\begin{lemma}[Weighted path cosines; proved here]
\label{lem:terminal-modal-cosines}
The vectors $u_h$ are an orthogonal eigenbasis of $P_{\rm row}$ with
eigenvalues $\mu_h=\cos(h\pi/m)$.  Their squared norms are
\begin{equation}
 N_h:=\norm{u_h}_{D}^{2}
 =\begin{cases}
 2m,&h\in\{0,m\},\\
 m,&1\le h\le m-1.
 \end{cases}
 \label{eq:terminal-modal-norms}
\end{equation}
Consequently every vector $z$ has coefficients
$\widehat z_h=\langle z,u_h\rangle_D/N_h$, and
\begin{equation}
 \norm{z-\widehat z_0\one}_D^2
 =m\sum_{h=1}^{m-1}\widehat z_h^2+2m\widehat z_m^2.
 \label{eq:terminal-modal-parseval}
\end{equation}
\end{lemma}

\begin{proof}
The endpoint equations are
$(P_{\rm row}u_h)_0=u_h(1)=\mu_hu_h(0)$ and similarly at $m$.
At an interior point, the cosine addition formula gives
$(u_h(j-1)+u_h(j+1))/2=\mu_hu_h(j)$.  Reversibility,
$DP_{\rm row}=A=A^\top$, makes distinct modes $D$-orthogonal.  Finally,
the standard discrete-cosine sums with endpoint weights one and interior
weights two give \eqref{eq:terminal-modal-norms}; Parseval follows.
\end{proof}

Let $a_{h,k}=\langle r_k^x,u_h\rangle_D/N_h$ and set
\begin{equation}
 \phi_h=\frac{h\pi}{2m},\qquad
 \varrho_h=(1-q)\cos\phi_h.
 \label{eq:terminal-modal-phase-radius}
\end{equation}

\begin{lemma}[Literal modal solution; proved here]
\label{lem:terminal-modal-solution}
For $1\le h\le m-1$, every projection-inactive full-face trajectory obeys
\begin{equation}
 a_{h,k}=\varrho_h^k
 \left(C_h\cos(k\phi_h)+D_h\sin(k\phi_h)\right),
 \label{eq:terminal-modal-solution}
\end{equation}
where
\begin{equation}
 C_h=a_{h,0},\qquad
 D_h=\frac{a_{h,1}/\varrho_h-C_h\cos\phi_h}{\sin\phi_h}.
 \label{eq:terminal-modal-quadratures}
\end{equation}
The constant mode has a repeated root $1-q$ and is irrelevant to residual
range.  The $h=m$ mode is killed after one application of $B$.
\end{lemma}

\begin{proof}
On mode $h$, the eigenvalue of $B$ is
$(1-q^2)\cos^2\phi_h$.  Thus
\cref{eq:terminal-modal-second-order} has characteristic polynomial
\[
 z^2-2(1-q)\cos^2\phi_h\,z
 +(1-q)^2\cos^2\phi_h,
\]
whose roots are $\varrho_he^{\pm i\phi_h}$.  Matching $k=0,1$ gives
\eqref{eq:terminal-modal-quadratures}.  The two endpoint cases follow by
direct substitution.
\end{proof}

The safe-envelope correction on the full face is
\begin{equation}
 \delta_k=\max\left(0,\frac{\max_j r_{k,j}^x}{\alpha}\right),
 \qquad
 \ell_k=[p_k-\delta_k\one]_+.
 \label{eq:terminal-modal-correction}
\end{equation}
This maximum is global; it is not a seed or frontier surrogate.

\begin{lemma}[Literal range implication; proved here]
\label{lem:terminal-modal-range}
Suppose $p_k-\delta_k\one>0$ coordinatewise.  Put
$R_k=\max_jr_{k,j}^x-\min_jr_{k,j}^x$.  If
\begin{equation}
 R_k>\alpha\tau=\frac{q^3}{5},
 \label{eq:terminal-modal-range-threshold}
\end{equation}
then the one-sided certificate fails at step $k$.
\end{lemma}

\begin{proof}
Because $P_{\rm row}\one=\one$, one has $\widehat Q\one=\alpha\one$.
Unclipping therefore gives
\[
 r(\ell_k)=r_k^x-M_k\one,
 \qquad M_k=\max(0,\max_jr_{k,j}^x).
\]
If $\max_jr_{k,j}^x\ge0$, the minimum corrected residual is exactly $-R_k$.
If the maximum is negative, the correction is zero and
$\min_jr_{k,j}^x< -R_k$.  In either case
\eqref{eq:terminal-modal-range-threshold} leaves an active row below
$-\alpha\tau$.
\end{proof}

Let $\bar r_k=\langle r_k^x,\one\rangle_D/(2m)$ and
$s_k=r_k^x-\bar r_k\one$.  Weighted Popoviciu gives
\begin{equation}
 R_k\ge\frac{2\norm{s_k}_D}{\sqrt{2m}}.
 \label{eq:terminal-modal-popoviciu}
\end{equation}
Indeed, the degree weights divided by $2m$ define a probability distribution,
and the variance of a variable in an interval of length $R_k$ is at most
$R_k^2/4$.  This is why modal $\ell_2$ energy prevents spatial cancellation
without selecting a particular vertex.

\section{The binomial packet and phase-uniform anti-cancellation}
\label{sec:terminal-modal-packet}

Define the ideal entry packet
\begin{equation}
 g_j=-\frac{\alpha}{2}(1-q)^m b_{m,j},
 \qquad b_{m,j}=2^{-m}\binom mj.
 \label{eq:terminal-modal-packet}
\end{equation}
Its center height is
\begin{equation}
 \frac{\max_j|g_j|}{q^{5/2}}
 \longrightarrow 2e^{-1/16}\sqrt{\frac2\pi}.
 \label{eq:terminal-modal-packet-height}
\end{equation}

\begin{lemma}[Exact packet coefficients; proved here]
\label{lem:terminal-modal-packet-transform}
For $1\le h\le m-1$, the weighted cosine coefficient of $g$ is
\begin{equation}
 \widehat g_h=-\frac{\alpha(1-q)^m}{m}
 \left[
  \cos^m\phi_h\cos\frac{h\pi}{2}
  -2^{-m-1}(1+(-1)^h)
 \right].
 \label{eq:terminal-modal-packet-transform}
\end{equation}
In particular, for $h=2s$,
\begin{equation}
 G_{2s}:=\widehat g_{2s}
 =-\frac{\alpha(1-q)^m}{m}
 \left[(-1)^s\cos^m\frac{s\pi}{m}-2^{-m}\right].
 \label{eq:terminal-modal-even-packet}
\end{equation}
For $s=o(\sqrt m)$,
$|G_{2s}|\sim e^{-1/16}\alpha/m=16e^{-1/16}q^3$.
\end{lemma}

\begin{proof}
If $J\sim\operatorname{Binomial}(m,1/2)$, then
\[
 \mathbb E e^{itJ}=e^{imt/2}\cos^m(t/2),
 \qquad
 \mathbb E\cos(tJ)=\cos(mt/2)\cos^m(t/2).
\]
The degree weights double every interior binomial mass and leave the two
endpoint masses single.  Taking $t=h\pi/m$, subtracting those endpoint
copies, multiplying by $-\alpha(1-q)^m/2$, and dividing by $N_h=m$ gives
\eqref{eq:terminal-modal-packet-transform}.  The even-mode specialization and
asymptotic follow immediately.
\end{proof}

\begin{lemma}[Uniform cosine-square floor; proved here]
\label{lem:terminal-modal-cosine-square}
For every integer $H\ge8$ and every real $x$,
\begin{equation}
 \sum_{s=1}^{H}\cos^2(sx)\ge\frac{H}{16}.
 \label{eq:terminal-modal-cosine-square}
\end{equation}
\end{lemma}

\begin{proof}
If $\sin x=0$, the sum equals $H$.  Otherwise, the finite sum identity gives
\[
 \sum_{s=1}^H\cos^2(sx)
 =\frac H2+\frac{\sin(Hx)\cos((H+1)x)}{2\sin x}.
\]
If $|\sin x|\ge2/H$, the right side is at least $H/4$.  Otherwise let
$\xi$ be the distance from $x$ to $\pi\mathbb Z$.  Since
$\sin(\pi/H)\ge2/H$, one has $\xi<\pi/H$.  For
$1\le s\le\lfloor H/4\rfloor$, $s\xi\le\pi/4$, so each corresponding
square is at least $1/2$.  Since
$\lfloor H/4\rfloor\ge H/8$ for $H\ge8$, the result follows.
\end{proof}

The next two statements are the exact open interfaces.  They are displayed as
lemmas to freeze the needed inequalities, but are not asserted as proved.

\begin{lemma}[Entry packet quadratures; open]
\label{lem:terminal-modal-packet}
Put
\begin{equation}
 H_m=\left\lfloor\sqrt{\frac{m}{64\log(16m)}}\right\rfloor.
 \label{eq:terminal-modal-bandwidth}
\end{equation}
For all sufficiently large $m$, the literal full-face entry coefficients from
\cref{eq:terminal-modal-quadratures} satisfy, for $1\le s\le H_m$,
\begin{equation}
 |C_{2s}-G_{2s}|+|D_{2s}|\le\frac{|G_{2s}|}{64}.
 \label{eq:terminal-modal-entry-target}
\end{equation}
\end{lemma}

The sine denominator in $D_{2s}$ is essential.  A bound only on
$r_0^x-g$ does not imply \eqref{eq:terminal-modal-entry-target}; one must also
control the entry velocity to order $\sin(s\pi/m)$.

\begin{lemma}[Projection and global-envelope regime; open]
\label{lem:terminal-modal-regime}
Put
\begin{equation}
 K_m=\left\lfloor\frac1{8q}\log\frac1q\right\rfloor.
 \label{eq:terminal-modal-horizon}
\end{equation}
For every $1\le k\le K_m$, the raw point
$y_{k-1}-r(y_{k-1})$ is strictly positive, and the resulting full-face state
satisfies
\begin{equation}
 \min_j p_{k,j}>\delta_k.
 \label{eq:terminal-modal-unclipped-target}
\end{equation}
\end{lemma}

This statement preserves both nonlinearities that the linear recurrence omits:
the coordinate projection and the moving global correction.

\begin{theorem}[Conditional terminal logarithmic block; proved here]
\label{thm:terminal-modal-conditional-log}
Assume \cref{lem:terminal-modal-packet,lem:terminal-modal-regime}.  Then, for
all sufficiently large $m$, the named execution has no terminal one-sided
certificate during its first $K_m$ full-face steps.  Hence its terminal block
has length
\begin{equation}
 \Omega\!\left(q^{-1}\log\frac1q\right).
 \label{eq:terminal-modal-conditional-log}
\end{equation}
This conclusion is for the named literal recurrence only.
\end{theorem}

\begin{proof}
For $1\le s\le H_m$, let $A_s=|G_{2s}|$.  Uniformly on this band and for all
sufficiently large $m$,
\begin{equation}
 \frac{\alpha}{3m}\le A_s\le\frac{2\alpha}{m}.
 \label{eq:terminal-modal-amplitude-bounds}
\end{equation}
Indeed, $(1-q)^m\ge1-mq=15/16$, while
$m(s\pi/m)^2\le\pi^2/(64\log(16m))$ and
$\log\cos z\ge-z^2$ for sufficiently small $z$.

Fix $1\le k\le K_m$.  For $h=2s$, the phase in
\cref{eq:terminal-modal-solution} is $\phi_{2s}=s\pi/m$.  Moreover,
\begin{equation}
 \cos^{k}\frac{H_m\pi}{m}
 \ge \exp\!\left(-\frac{\pi^2}{32}\right)>\frac12.
 \label{eq:terminal-modal-spatial-damping}
\end{equation}
Apply \cref{lem:terminal-modal-cosine-square} at
$x=k\pi/m$.  The ideal part of the band coefficient vector has Euclidean norm
at least
\[
 (1-q)^k\cdot\frac{\alpha}{3m}\cdot\frac12
 \sqrt{\frac{H_m}{16}}
 =\frac{\alpha(1-q)^k\sqrt{H_m}}{24m}.
\]
By \eqref{eq:terminal-modal-entry-target}, its perturbation has norm at most
\[
 (1-q)^k\cdot\frac{2\alpha}{m}\cdot
 \frac{\sqrt{H_m}}{64}
 =\frac{\alpha(1-q)^k\sqrt{H_m}}{32m}.
\]
The reverse triangle inequality therefore gives
\begin{equation}
 \left(\sum_{s=1}^{H_m}a_{2s,k}^2\right)^{1/2}
 \ge\frac{\alpha(1-q)^k\sqrt{H_m}}{96m}.
 \label{eq:terminal-modal-band-energy}
\end{equation}
Parseval and \cref{eq:terminal-modal-popoviciu} imply
\begin{equation}
 R_k\ge
 \frac{\sqrt2\alpha(1-q)^k\sqrt{H_m}}{96m}
 =\frac{\sqrt2}{6}q^3(1-q)^k\sqrt{H_m}.
 \label{eq:terminal-modal-range-lower}
\end{equation}

Because $q\le1/32$,
\[
 (1-q)^{K_m}
 \ge \exp\!\left(-\frac{qK_m}{1-q}\right)
 \ge q^{4/31}.
\]
Also
$\sqrt{H_m}=\Omega(q^{-1/4}(\log(1/q))^{-1/4})$.  Consequently the ratio of
the right side of \eqref{eq:terminal-modal-range-lower} to $q^3$ diverges
like a positive constant times
$q^{-15/124}(\log(1/q))^{-1/4}$.  It is therefore greater than $1/5$ for all
sufficiently small $q$.  The open regime lemma permits
\cref{lem:terminal-modal-range}, which proves non-certification through
$K_m$.
\end{proof}

\begin{remark}[Why one mode is insufficient]
For each fixed $h$, the ideal coefficient is only $\Theta(q^3)$, exactly the
certificate scale, and it can pass through zero.  The gain comes from a band
of $H_m$ orthogonal modes and the phase-uniform
\cref{lem:terminal-modal-cosine-square}; it is not a surviving-mode argument.
The conservative proof constant $1/8$ is not intended to predict the measured
first-certificate constant, which is closer to
$\frac12\log(1/q)$ over the screened sizes.
\end{remark}

\section{Deterministic screen}
\label{sec:terminal-modal-screen}

The note-local script \texttt{verify.py} reconstructs every singleton
admission using a Thomas solve for each transported restricted optimum, then
runs the literal projected full-face recurrence.  It records both the first
global crossing $R_k\le q^3/5$ and the first literal corrected test
$\min r_k-\max(0,\max r_k)\ge-q^3/5$.  It separately checks the minimum raw
projection and $p_k-\delta_k\one$ margins.  It uses deterministic float64
arithmetic and therefore supplies measurements, not exact or asymptotic proof.
On the exact band $1\le s\le H_m$, it also reports and checks the combined
quantity
\[
 \max_s\frac{|C_{2s}-G_{2s}|+|D_{2s}|}{|G_{2s}|}\le\frac1{64},
\]
with a vacuous label when $H_m=0$.  The separately displayed ``modal defect''
uses the wider diagnostic band described below and is not that direct check.

\begin{center}
\begin{tabular}{@{}r|rrrrrrr@{}}
$m$ & $\norm{r_0}_{1,D}/q^2$ & $R_0/q^{5/2}$
& $\norm{r_0-g}_{1,D}/q^2$ & modal defect & $qk_{\rm range}$
& $qk_{\rm cert}$
& margin$/q^2$\\
\hline
128  & .972971 & 1.497735 & .033572 & .00814 & 3.123535 & 3.733887 & 7.7761\\
256  & .972979 & 1.498840 & .033573 & .00808 & 3.871094 & 3.985840 & 7.7776\\
512  & .972983 & 1.499267 & .033574 & .00805 & 4.366577 & 4.366577 & 7.7784\\
1024 & .972985 & 1.499392 & .033574 & .00803 & 4.713074 & 4.713074 & 7.7788\\
2048 & .972986 & 1.499391 & .033574 & .00802 & 5.017151 & 5.017151 & 7.7790\\
\end{tabular}
\end{center}

Here ``modal defect'' is the maximum of
$|C_{2s}-G_{2s}|/|G_{2s}|$ and $|D_{2s}|/|G_{2s}|$ over
$1\le s\le\lfloor\sqrt{m/\log m}\rfloor$, a wider band than the conditional
theorem needs.  The final margin is the smaller of the raw-projection and
unclipped-envelope margins through the literal certificate.  For $m=128,256$,
the first range crossing has a negative residual maximum and the minimum
remains below $-q^3/5$; it is therefore not a literal certificate.  For the
other displayed sizes, the two events coincide.  This distinction is why the
script evaluates both tests rather than silently identifying them.

The entry constants are explained by the ideal packet:
$\norm{g}_{1,D}/q^2\to e^{-1/16}$ and
$\operatorname{range}(g)/q^{5/2}\to
2e^{-1/16}\sqrt{2/\pi}$.  The measured residual has a small, spatially smooth
remainder.  This is evidence for
\cref{lem:terminal-modal-packet}, not a bound on that remainder's entry
velocity.

\section{Falsifiable conclusion and next proof step}
\label{sec:terminal-modal-conclusion}

The root calculation is no longer the blocker.  Exact orthogonality removes
spatial cancellation in weighted $\ell_2$, and a consecutive even-mode band
removes temporal phase cancellation.  A theorem-quality logarithmic block now
reduces to precisely two uniform statements:
\begin{enumerate}
 \item derive \eqref{eq:terminal-modal-entry-target} from the literal
 admission/transport chronology, including the sine-scaled entry velocity;
 \item prove \eqref{eq:terminal-modal-unclipped-target} and raw positivity
 through \eqref{eq:terminal-modal-horizon}.
\end{enumerate}

Failure of either displayed inequality for an unbounded subsequence of $m$
falsifies this route.  Success proves only a lower block for the named
transported-center, full-prefix, exact-real recurrence.  It would not imply a
lower bound for distributed residual mechanisms, response data structures,
different projection rules, intermediate external returns, finite precision,
or the broad class of local accelerated algorithms.

\end{document}
